%% Determines the type of document (including standard settings for layout).
\documentclass[letterpaper]{article}



%% Package to control the size of a page: the standard margins used by LaTeX are
%% very wide!
\usepackage[margin=1in]{geometry}


%% Packages add functionality to the document. The AMS packages are standard
%% packages to support various mathematical notations. AMS stands for American
%% Mathematical Society, the organization that maintains these packages.
\usepackage{amsmath,amsthm,amssymb}

%% Theorems-like environments using functionality provided by amsthm.
\theoremstyle{plain}
\newtheorem{theorem}{Theorem}[section]
    %% [section] at the end specifies that theorems should be numbered
    %% per-section: Section x starts with theorem-like x.1 and so on, ...
\newtheorem{proposition}[theorem]{Proposition}
    %% [theorem] in the middle specifies: use the same counter as the theorem
    %% environment: here we number all theorem-like environments consecutively.
\newtheorem{corollary}[theorem]{Corollary}
\newtheorem{lemma}[theorem]{Lemma}
\theoremstyle{definition}
\newtheorem{definition}[theorem]{Definition}
\theoremstyle{remark}
\newtheorem{example}[theorem]{Example}
\newtheorem{remark}[theorem]{Remark}


%% Support for nicely formatted tables.
\usepackage{booktabs}


%% Support for colors & colors in tables.
\usepackage[table]{xcolor}

% Seven colors safe for use color blindness.
% Colors taken from doi:10.1038/nmeth.1618.
\definecolor{cbOrange}{RGB}{230,159,0}
\definecolor{cbSkyBlue}{RGB}{86,180,233}
\definecolor{cbBluishGreen}{RGB}{0,158,115}
\definecolor{cbYellow}{RGB}{240,228,66}
\definecolor{cbBlue}{RGB}{0,114,178}
\definecolor{cbVermillion}{RGB}{213,94,0}
\definecolor{cbReddischPurple}{RGB}{204,121,167}


%% Notation used in this document.
\newcommand{\n}{\mathbf{n}} %% Num. Replicas.
\newcommand{\f}{\mathbf{f}} %% Num. Faulty Replicas.

%% Misc. Math notation.
\newcommand{\BigO}{\mathcal{O}}
\newcommand{\abs}[1]{\lvert #1 \rvert}
\newcommand{\AName}[1]{\textsc{#1}}
\newcommand{\Var}[1]{\texttt{#1}}


%% Algorithms.
\usepackage{algorithm}
\usepackage[noend]{algorithmic}
\newcommand{\GETS}{:=}


%% Formatting SI-units.
\usepackage{siunitx}
\sisetup{per-mode=symbol}


%% TikZ: for creating figures.
\usepackage{tikz}

%% Configuration for figures: Nicer arrows.
\usetikzlibrary{arrows.meta}
\tikzset{>=Stealth}


%% pgfplots: drawing plots using TikZ.
\usepackage{pgfplots}
%% Configuration for plots: Use color-blind friendly colors.
\pgfplotscreateplotcyclelist{cbSafeList}{
    very thick,solid,cbOrange,every mark/.append style={solid},mark=*\\
    very thick,solid,cbSkyBlue,every mark/.append style={solid},mark=*\\
    very thick,solid,cbBluishGreen,every mark/.append style={solid},mark=*\\
    very thick,solid,cbYellow,every mark/.append style={solid},mark=*\\
    very thick,solid,cbBlue,every mark/.append style={solid},mark=*\\
    very thick,solid,cbVermillion,every mark/.append style={solid},mark=*\\
    very thick,solid,cbReddischPurple,every mark/.append style={solid},mark=*\\
    very thick,solid,black,every mark/.append style={solid},mark=*\\
}
\pgfplotsset{
    legend style={font=\small},
    compat=1.16,
    width=260pt,
    height=140pt,
    legend cell align=left,
    xlabel near ticks,
    ylabel near ticks,
    every axis/.append style={
        cycle list name=cbSafeList,
        ymin=0,
        enlargelimits=0.05,
        mark size=1pt,
        ylabel style={align=center},
        xlabel style={align=center},
        title style={align=center}
    }
}

%% PgfplotsTable: loading data files to use with pgfplots.
\usepackage{pgfplotstable}


%% Support for hyperlinks and urls. The setting ``colorlinks'' sets how links
%% are shown in the document (with a color, without underline). We put hyperref
%% last---it has a tendency to break other packages when loaded before them.
\usepackage[colorlinks]{hyperref}
\usepackage{graphicx}
\usepackage{pgffor}
\usepackage{caption}
\usepackage{tabularx}
\usepackage{enumitem}
\usepackage{fancyhdr}

\title{COMPSCI 2AC3 Assignment 4}
\author{Luca Mawyin - 400531739}
\date{\today}

\begin{document}

\maketitle

Chapters to review: 22, 23, 24, 25, 27, 28, 29, 31
\renewcommand{\thesubsection}{\thesection.\alph{subsection}}

\stepcounter{section}
\stepcounter{subsection}
\section*{\thesubsection}

\begin{align*}
    A&=\{b^ka^{2k}b^{3k}~|~k\geq0\}
\end{align*}
The given set is not a CFL. We will use the pumping lemma and demon game to prove.
\begin{enumerate}
    \item Suppose the Demon chooses any $k\geq0$
    \item We will pick $z=b^ka^{2k}b^{3k}$, as $z\in{A}\land|z|=6k\geq{k}$, trivially
    \item Suppose the Demon picks some $u,v,w,x,y$ such that $z=uvwxy\land vx\neq\epsilon\land|vwx|\leq{k}$
    \item We will choose  $i=2$. Then:
    \begin{align*}
    z&=uv^2wx^2y
    \end{align*}
    In every case, this cannot be expressed in the form $b^ka^{2k}b^{3k}$.
\end{enumerate}
As we cannot express the set as shown by the pumping lemma, the set is not a CFL. 

\stepcounter{subsection}
\section*{\thesubsection}

\begin{align*}
    B&=\{a^nb^m~|~n\geq3\text{ and }m\geq{n}\}
\end{align*}
The given set is a CFL. We can define the grammar as $G=(N,\Sigma,P,S)$, where:
\begin{align*}
    N&=\{S,T\}\\
    \Sigma&=\{a,b\}\\
    P&=
    \left\{
    \begin{array}{l}
        S \to aaaTbbb\\
        T \to aTb~|~Tb~|~\epsilon
    \end{array}
    \right.
\end{align*}

\section{}
\begin{align*}
    C&=\{a^mb^nc^kd^m~|~n\geq0\text{ is even and } m,k\geq1\}
\end{align*}
As every CFG is accepted by some NPDA, we can first define the grammar for the language as \\
$G=(N,\Sigma,P,S)$, where:
\begin{align*}
    N&=\{S,T,U\}\\
    \Sigma&=\{a,b,c,d\}\\
    P&=
    \left\{
    \begin{array}{l}
        S \to aSd~|~aTd \\
        T \to bbT~|~U\\
        U \to c~|~cU
    \end{array}
    \right.
\end{align*}
\newpage
We will let $M=(Q,\Sigma,\Gamma,\delta,s,\bot,F)$ be our NPDA where:
\begin{align*}
    Q&=\{q_s,q,q_f\}\\
    \Sigma&=\{a,b,c,d\}\\
    \Gamma&=\{S,T,U,V,\bot\}\\
    s&=q_s\\
    F&=q_f\\
\end{align*}
Now we will define our transition functions $\delta$:
\begin{itemize}
    \item $\delta(q_s,\epsilon,\bot)=(q,S\bot)$
    \item $S\to aSd$ is $\delta(q,a,S)=(q,Sd)$
    \item $S\to aTd$ is $\delta(q,a,S)=(q,Td)$
    \item $T\to{bbT}$ is split into $T\to{bV}$ and $V\to{bT}$
    \item $T\to{bV}$ is $\delta(q,b,T)=(q,V)$
    \item $V\to{bT}$ is $\delta(q,b,V)=(q,T)$
    \item $T\to{U}$ is $\delta(q,\epsilon,T)=(q,U)$
    \item $U\to{cU}$ is $\delta(q,c,U)=(q,U)$
    \item $U\to{c}$ is $\delta(q,c,U)=(q,\epsilon)$
    \item $\delta(q,\epsilon,\bot)=(q_f,\epsilon)$
\end{itemize}
For the NPDA, there are three states $q_s$, $q$, and $q_f$. This is because, without there being $q_s$, it would be possible for a string to be empty and pass the NPDA, which goes against the constrains of our language. Furthermore, matching the transitions of the CFG to those of the NPDA, the transitions are extremely intuitive. The stack reads the character, pushes the corresponding transition to the stack, and proceeds to the next character. Once the stack is complete (which coincides with completion of the string), The NPDA makes an $\epsilon$-transition to the accepting state.

\section{}
\begin{align*}
    &S \to bD_2~|~cD_3~|~\epsilon\\
    &D_2 \to aa~|~D_2aa~|~\epsilon\\
    &D_3 \to aaa~|~D_3aaa~|~\epsilon
\end{align*}

\stepcounter{subsection}
\section*{\thesubsection}
The grammar creates string generates strings in the form $ba^*$ where $\#a\mod 2 = 0$, or $ca^*$ where $\#a\mod 3 = 0$. More formally, we can say:

\begin{align*}
    L(G)&=\{ba^n~|~n\text{ mod }2\equiv0\}~\cup~\{ca^m~|~m\text{ mod }3\equiv0\}~\cup~\{\epsilon\}
\end{align*}

\stepcounter{subsection}
\section*{\thesubsection}
\begin{enumerate}
    \item Add nonterminals and remove $\epsilon$
    \begin{align*}
        &S \to BD_2~|~CD_3\\
        &D_2 \to AA~|~D_2AA\\
        &D_3 \to AAA~|~D_3AAA\\
        &A \to a\\
        &B \to b\\
        &C \to c
    \end{align*} 
    \item Add nonterminal $T$,$U$, and $V$ and replace $D_2$ and $D_3$ with:
    \begin{align*}
        &S \to BD_2~|~CD_3\\
        &D_2 \to U~|~TA\\
        &D_3 \to AU~|~VU\\
        &V \to D_3A\\
        &T \to D_2A\\
        &U \to AA\\
        &A \to a\\
        &B \to b\\
        &C \to c
    \end{align*} 
\end{enumerate}

\stepcounter{subsection}
\section*{\thesubsection}

$x = $ \begin{tabular}{|l|l|l|l|l|l|l|}
    b&a&a&a&a&a&a
\end{tabular}\\[1em]
We have $n=|x|=7$, therefore we hav our table as follows:\\[1em]
\begin{tabular}{@{}c c c c c c c c@{}}
0 &   &   &   &   &   &   &   \\
- & 1 &   &   &   &   &   &   \\
- & - & 2 &   &   &   &   &   \\
- & - & - & 3 &   &   &   &   \\
- & - & - & - & 4 &   &   &   \\
- & - & - & - & - & 5 &   &   \\
- & - & - & - & - & - & 6 &   \\
- & - & - & - & - & - & - & 7 \\
\end{tabular}\\[1em]
Starting with substrings of length 1:\\[1em]
\begin{tabular}{@{}c c c c c c c c@{}}
0 &   &   &   &   &   &   &   \\
B & 1 &   &   &   &   &   &   \\
- & A & 2 &   &   &   &   &   \\
- & - & A & 3 &   &   &   &   \\
- & - & - & A & 4 &   &   &   \\
- & - & - & - & A & 5 &   &   \\
- & - & - & - & - & A & 6 &   \\
- & - & - & - & - & - & A & 7 \\
\end{tabular}\\[5em]
Substrings of length 2:\\[1em]
\begin{tabular}{@{}c c c c c c c c@{}}
0 &   &   &   &   &   &   &   \\
B & 1 &   &   &   &   &   &   \\
$\varnothing$ & A & 2 &   &   &   &   &   \\
- & $\{D_2,U\}$& A & 3 &   &   &   &   \\
- & - & $\{D_2,U\}$ & A & 4 &   &   &   \\
- & - & - & $\{D_2,U\}$ & A & 5 &   &   \\
- & - & - & - & $\{D_2,U\}$ & A & 6 &   \\
- & - & - & - & - & $\{D_2,U\}$ & A & 7 \\
\end{tabular}\\[1em]
Substrings of length 3:\\[1em]
\begin{tabular}{@{}c c c c c c c c@{}}
0 &   &   &   &   &   &   &   \\
B & 1 &   &   &   &   &   &   \\
$\varnothing$ & A & 2 &   &   &   &   &   \\
$S$ & $\{D_2,U\}$& A & 3 &   &   &   &   \\
- & $\{D_3,T\}$ & $\{D_2,U\}$ & A & 4 &   &   &   \\
- & - & $\{D_3,T\}$ & $\{D_2,U\}$ & A & 5 &   &   \\
- & - & - & $\{D_3,T\}$ & $\{D_2,U\}$ & A & 6 &   \\
- & - & - & - & $\{D_3,T\}$ & $\{D_2,U\}$ & A & 7 \\
\end{tabular}\\[1em]
Substrings of length 4:\\[1em]
\begin{tabular}{@{}c c c c c c c c@{}}
0 &   &   &   &   &   &   &   \\
B & 1 &   &   &   &   &   &   \\
$\varnothing$ & A & 2 &   &   &   &   &   \\
$S$ & $\{D_2,U\}$& A & 3 &   &   &   &   \\
$\varnothing$ & $\{D_3,T\}$ & $\{D_2,U\}$ & A & 4 &   &   &   \\
$S$ & $D_2$ & $\{D_3,T\}$ & $\{D_2,U\}$ & A & 5 &   &   \\
$\varnothing$ & $\varnothing$ & $D_2$ & $\{D_3,T\}$ & $\{D_2,U\}$ & A & 6 &   \\
$S$ & $D_3$ & $\varnothing$ & $D_2$ & $\{D_3,T\}$ & $\{D_2,U\}$ & A & 7 \\
\end{tabular}\\[1em]
Based on the table drawn, as $T(0,7)$ contains the start symbol $S$, this means that the string $x$ can be generated by $G'$

\section{}
\begin{center}
    $D=\{a^nb^mc^{n\cdot{m}}~|~n,m\geq1\}$    
\end{center}
We will define our Turing maching as $M=(Q,\Sigma,\Gamma,\vdash,\sqcup,\delta,s,t,r)$ where:
\begin{align*}
    Q&=\{s,q_1,\ldots,q_{10},t,r\}\\
    \Sigma&=\{a,b,c\}\\
    \Gamma&=\Sigma\cup\{X,Y,Z,\vdash,\sqcup,\dashv\}
\end{align*}
Now we will create our transition table\\[1em]
\begin{tabular}{c | l l l l l l l l l}
    & $\vdash$ & $a$ & $b$ & $c$ & $X$ & $Y$ & $Z$ & $\sqcup$ & $\dashv$ \\
    \hline
\end{tabular}
\end{document}